% ============================================================
% РОЗДІЛ 3 — переписаний (структура без підпідрозділів, стиль академічний)
% Вставити замість блоку від \protect\phantomsection\label{_Toc229681122}{}
% до \protect\phantomsection\label{_Toc229681164}{} у БКР_Тесля.tex
% ============================================================

\protect\phantomsection\label{_Toc229681122}{}РОЗДІЛ 3. РЕАЛІЗАЦІЯ
ПРОГРАМНОГО РІШЕННЯ

\protect\phantomsection\label{_Toc229681123}{}3.1. Опис технологій та
інструментів розробки

У підрозд.~2.1 обґрунтовано архітектурні рішення: клієнтський SPA,
реактивний рендеринг, централізоване сховище стану, спеціалізована
бібліотека графової візуалізації. У цьому підрозділі наведено конкретні
технології, бібліотеки та їхні версії, що реалізують зазначені
архітектурні рішення.

Основною мовою реалізації обрано \textbf{TypeScript~5.x} --- надбудову
над JavaScript зі статичною типізацією. Мова компілюється у
стандартизований ECMAScript~2020 (ES11), який виконується усіма сучасними
браузерами без додаткових polyfill'ів (NF4). Статична типізація дозволяє
виявляти на етапі компіляції помилки передачі об'єкта \texttt{Token}
замість числа чи неправильної форми комплексного значення, що критично для
дотримання NF1 (\(\varepsilon\leq10^{-9}\)). Цільовим середовищем
виконання є браузерний рушій JavaScript (V8 у Chrome/Edge, SpiderMonkey у
Firefox) з підтримкою \texttt{requestAnimationFrame}, Web Workers,
Canvas API та Local Storage.

Повний перелік використаних бібліотек з версіями та призначенням наведено
у табл.~3.1. \textbf{Next.js} забезпечує маршрутизацію, оптимізоване
завантаження ресурсів і статичну збірку без серверної частини (режим
\texttt{output: 'export'}); усі сторінки позначені директивою
\texttt{'use client'}, оскільки \texttt{requestAnimationFrame}-цикл
несумісний зі server-side rendering. \textbf{React~Flow} надає готові
механізми перетягування вузлів, масштабування, анімованих ребер і
кастомних типів вузлів; його модель \texttt{Node}/\texttt{Edge}
безпосередньо відображається на формальну модель DFG (підрозд.~2.4).
\textbf{Zustand} --- легка альтернатива Redux без шаблонного коду
редьюсерів; компоненти підписуються через хук
\texttt{useSimStore(selector)} з мемоїзацією. \textbf{fft.js}
використовується виключно як референтна реалізація ДПФ для верифікації і
не залучається до шляху виконання основного алгоритму.

Таблиця 3.1.

Технологічний стек реалізації

\begin{longtable}[]{@{}
  >{\raggedright\arraybackslash}p{(\linewidth - 4\tabcolsep) * \real{0.3293}}
  >{\raggedright\arraybackslash}p{(\linewidth - 4\tabcolsep) * \real{0.1068}}
  >{\raggedright\arraybackslash}p{(\linewidth - 4\tabcolsep) * \real{0.5639}}@{}}
\caption[Технологічний стек реалізації]{Технологічний стек
реалізації}\tabularnewline
\toprule\noalign{}
\begin{minipage}[b]{\linewidth}\raggedright
\textbf{Технологія}
\end{minipage} & \begin{minipage}[b]{\linewidth}\raggedright
\textbf{Версія}
\end{minipage} & \begin{minipage}[b]{\linewidth}\raggedright
\textbf{Призначення}
\end{minipage} \\
\midrule\noalign{}
\endfirsthead
\toprule\noalign{}
\begin{minipage}[b]{\linewidth}\raggedright
\textbf{Технологія}
\end{minipage} & \begin{minipage}[b]{\linewidth}\raggedright
\textbf{Версія}
\end{minipage} & \begin{minipage}[b]{\linewidth}\raggedright
\textbf{Призначення}
\end{minipage} \\
\midrule\noalign{}
\endhead
\bottomrule\noalign{}
\endlastfoot
TypeScript & 5.x & Статично типізована мова реалізації \\
Next.js & 14.x & React-фреймворк, статичний експорт SPA \\
React & 18.x & Реактивний рендеринг UI-компонентів \\
React Flow (\texttt{@xyflow/react}) & 12.x & Інтерактивна графова
візуалізація DFG \\
Zustand & 4.x & Централізоване реактивне сховище стану \\
Comlink & 4.x & RPC-обгортка для Web Workers \\
fft.js & 4.x & Референтна реалізація ДПФ для верифікації \\
Jest & 29.x & Фреймворк модульного тестування \\
ESLint & 8.x & Статичний аналіз коду, правила стилю \\
Prettier & 3.x & Автоматичне форматування коду \\
\end{longtable}

Розробка виконується у \textbf{Visual Studio Code} з розширеннями
ESLint, Prettier і TypeScript. Pre-commit-хуки реалізовано через
\texttt{husky}~+~\texttt{lint-staged}: перед кожним коммітом автоматично
запускаються ESLint та Prettier для змінених файлів. Менеджер пакетів ---
\textbf{pnpm}; файл \texttt{pnpm-lock.yaml} гарантує детермінованість
збірки на різних машинах. Використовується \textbf{Git} зі стратегією
trunk-based development: основна розробка ведеться у гілці \texttt{main},
для великих змін створюються короткоживучі feature-гілки. Репозиторій
розміщений на GitHub; опційно налаштовано \textbf{GitHub Actions} для
автоматичного запуску модульних тестів і лінтерів на кожному push до
\texttt{main} (підрозд.~3.4).

\protect\phantomsection\label{_Toc229681128}{}3.2. Структура проєкту та
організація вихідного коду

Структура кореневого каталогу проєкту організована за стандартом
Next.js з окремими директоріями для тестів і документації (рис.~3.1).

\begin{longtable}[]{@{}
  >{\raggedright\arraybackslash}p{(\linewidth - 0\tabcolsep) * \real{0.7080}}@{}}
\caption{Структура директорій проєкту. Джерело: розроблено
автором.}\tabularnewline
\toprule\noalign{}
\endfirsthead
\endhead
\bottomrule\noalign{}
\endlastfoot
\begin{minipage}[t]{\linewidth}\raggedright
\begin{verbatim}
simulator/
+-- public/                  # статичні ресурси (SVG-іконки, favicons)
+-- src/
|   +-- app/                 # Next.js App Router
|   |   +-- layout.tsx
|   |   \-- page.tsx
|   +-- core/                # ядро симуляції (незалежне від UI)
|   |   +-- types.ts
|   |   +-- scheduler.ts
|   |   \-- buffers.ts
|   +-- graph/               # генерація та візуалізація графа
|   |   +-- generateDFT4x4.ts
|   |   +-- layout4x4.ts
|   |   +-- GraphView.tsx
|   |   \-- TokenEdge.tsx
|   +-- components/          # UI-компоненти
|   |   +-- SpectrumView.tsx
|   |   +-- Metrics.tsx
|   |   +-- Controls.tsx
|   |   +-- TerminalOutput.tsx
|   |   \-- nodes/
|   |       +-- SourceNode.tsx
|   |       +-- SinkNode.tsx
|   |       +-- DFT4Node.tsx
|   |       \-- TwiddleNode.tsx
|   +-- store/               # централізоване сховище стану
|   |   \-- simStore.ts
|   +-- utils/               # допоміжні утиліти
|   |   +-- math.ts
|   |   \-- compare.ts
|   \-- workers/             # Web Workers
|       \-- fft.worker.ts
+-- tests/                   # Jest-тести
|   +-- math.test.ts
|   +-- scheduler.test.ts
|   \-- compare.test.ts
+-- package.json
+-- tsconfig.json
+-- next.config.js
\-- .eslintrc.js
\end{verbatim}
\end{minipage} \\
\end{longtable}

Рис.~3.1. Структура директорій проєкту. Джерело: розроблено автором.

Декомпозиція на модулі виконана за принципом \emph{розподілу за шарами
відповідальності}, що відповідає архітектурі Container з підрозд.~2.3.
Директорія \texttt{src/core/} містить ядро симуляції без будь-яких
залежностей від UI: воно може бути імпортоване у Node.js для
batch-обчислень і використовується в unit-тестах. Директорія
\texttt{src/graph/} охоплює генерацію структури DFG і компоненти
візуалізації на React Flow. Директорія \texttt{src/components/} містить
UI-компоненти, що не залежать від React Flow (панелі керування, метрики,
спектр). Сховище \texttt{src/store/} реалізує реактивний стан та
інтеграцію з ядром симуляції. Директорія \texttt{src/utils/} містить чисті
функції без побічних ефектів, а \texttt{src/workers/} --- Web Worker для
фонових обчислень. Така організація спрощує ізольоване тестування і
полегшує пошук місць для внесення змін.

Відповідність компонентів архітектури з підрозд.~2.3 та модулів
вихідного коду зведена у табл.~3.2.

Таблиця 3.2.

Відповідність «Компонент → Модуль/файл»

\begin{longtable}[]{@{}
  >{\raggedright\arraybackslash}p{(\linewidth - 4\tabcolsep) * \real{0.2555}}
  >{\raggedright\arraybackslash}p{(\linewidth - 4\tabcolsep) * \real{0.3745}}
  >{\raggedright\arraybackslash}p{(\linewidth - 4\tabcolsep) * \real{0.3700}}@{}}
\caption[Відповідність компонентів архітектури та модулів коду]{Відповідність
«Компонент~$\rightarrow$ Модуль/файл»}\tabularnewline
\toprule\noalign{}
\begin{minipage}[b]{\linewidth}\raggedright
\textbf{Компонент архітектури (2.3)}
\end{minipage} & \begin{minipage}[b]{\linewidth}\raggedright
\textbf{Модуль / файл}
\end{minipage} & \begin{minipage}[b]{\linewidth}\raggedright
\textbf{Роль}
\end{minipage} \\
\midrule\noalign{}
\endfirsthead
\toprule\noalign{}
\begin{minipage}[b]{\linewidth}\raggedright
\textbf{Компонент архітектури (2.3)}
\end{minipage} & \begin{minipage}[b]{\linewidth}\raggedright
\textbf{Модуль / файл}
\end{minipage} & \begin{minipage}[b]{\linewidth}\raggedright
\textbf{Роль}
\end{minipage} \\
\midrule\noalign{}
\endhead
\bottomrule\noalign{}
\endlastfoot
Simulation Core & \texttt{src/core/scheduler.ts} & Клас
\texttt{DataflowRuntime}, метод \texttt{tick()} \\
Simulation Core & \texttt{src/core/buffers.ts} & FIFO-черги на дугах \\
Simulation Core & \texttt{src/core/types.ts} & Типи \texttt{Token},
\texttt{Edge}, \texttt{NodeSpec} \\
Graph Generator & \texttt{src/graph/generateDFT4x4.ts} & Побудова DFG \\
Graph Generator & \texttt{src/graph/layout4x4.ts} & Координати вузлів \\
State Store & \texttt{src/store/simStore.ts} & Zustand-сховище \\
UI Layer (GraphView) & \texttt{src/graph/GraphView.tsx} & React Flow +
rAF-цикл \\
UI Layer (SpectrumView) & \texttt{src/components/SpectrumView.tsx} &
Canvas-діаграма \(|X(k)|\) \\
UI Layer (Metrics) & \texttt{src/components/Metrics.tsx} &
EMA-метрики \\
UI Layer (Controls) & \texttt{src/components/Controls.tsx} &
Пуск/пауза/швидкість \\
UI Layer (вузли) & \texttt{src/components/nodes/*.tsx} & SourceNode,
SinkNode, DFT4Node, TwiddleNode \\
Math \& verify & \texttt{src/utils/math.ts}, \texttt{compare.ts} &
Комплексна арифметика, верифікація \\
\end{longtable}

Конфігураційні файли визначають правила якості коду. Файл
\texttt{tsconfig.json} вмикає \texttt{strict}, \texttt{noImplicitAny},
\texttt{strictNullChecks} для максимальної безпеки типів; цільовий
стандарт~--- ES2020. Файл \texttt{next.config.js} встановлює
\texttt{output: 'export'} для статичної збірки та налаштовує базовий
шлях. Файл \texttt{.eslintrc.js} забороняє \texttt{any},
\texttt{console.log} у production, не імпортовані змінні та побічні
ефекти у хуках без залежностей. Ідентифікатори у коді --- англомовні,
камел-кейс для змінних і функцій, паскаль-кейс для типів і
React-компонентів; максимальна довжина рядка 100 символів контролюється
Prettier.

\protect\phantomsection\label{_Toc229681134}{}3.3. Реалізація ключових
модулів та алгоритмів

Клас \texttt{DataflowRuntime} (модуль \texttt{src/core/scheduler.ts}) є
програмною реалізацією абстрактної машини потоку даних із підрозд.~2.2
(Veen 1986; Dennis and Misunas 1975). Він відповідає за виконання правила
активації вузлів та управління FIFO-чергами на дугах графа і зберігає
такі внутрішні структури: незмінне посилання на граф \texttt{Graph};
відображення \texttt{Map<edgeId, FIFO<Token>>}~--- окрема черга токенів
для кожної дуги; \texttt{readyAt: Map<nodeId, number>}~--- час, раніше
якого вузол не може активуватись повторно (моделює \texttt{latency});
\texttt{inbox: Map<nodeId, Map<portId, Token>>}~--- токени, що вже
доставлені до вузла, але ще не поглинуті; обробники подій
\texttt{onToken}, \texttt{onFire}, \texttt{onBeforeFire},
\texttt{onOutput}. Логічний час надається ззовні через функцію
\texttt{getTime: () => number}, переданої у конструктор; у симуляторі нею
є \texttt{() => store.simTime}, завдяки чому годинник ядра повністю
відокремлений від астрономічного часу.

Метод \texttt{tick(dt, maxFires)} виконує єдиний крок потокового
обчислення і реалізує правило активації (Veen 1986). Алгоритм складається
з чотирьох послідовних кроків. По-перше, \textbf{доставка токенів}: для
кожної дуги \(e\) перевіряється умова
\(t_{now}-token.t\geq e.delay\); якщо виконана, токен видаляється з черги
і поміщається у \texttt{inbox} вузла-приймача, емітується подія
\texttt{onToken}. По-друге, \textbf{перевірка готовності}: вузол є
готовим, якщо у \texttt{inbox}~є токен на кожному вхідному порті та
\(\text{readyAt}[n]\leq t_{now}\). По-третє, \textbf{активація}: для
кожного готового вузла (не більше \texttt{maxFires} за виклик) токени
поглинаються з \texttt{inbox}, виконується логіка вузла за типом
\texttt{NodeKind}, час блокування оновлюється:
\(\text{readyAt}[n]\leftarrow t_{now}+n.latency\), емітуються події
\texttt{onBeforeFire} і \texttt{onFire}. По-четверте,
\textbf{розміщення вихідних токенів}: для кожного вихідного токена
знаходяться всі дуги з відповідним портом-джерелом і токен поміщається у
їхні черги з часом народження \(token.t = t_{now}+n.latency\), що
гарантує правильну затримку доставки. Такий порядок кроків забезпечує
детермінованість (NF5): при фіксованому обході вузлів результат симуляції
не залежить від реального часу у браузері.

Клас \texttt{FIFO<T>} (модуль \texttt{src/core/buffers.ts}) --- узагальнена
черга на базі масиву з методами \texttt{push(item)},
\texttt{pop(): T|undefined}, \texttt{peek(): T|undefined} та властивістю
\texttt{size}. Одна черга ставиться у відповідність кожній дузі графа і
ініціалізується у конструкторі \texttt{DataflowRuntime} ітерацією по
\texttt{graph.edges}. При характерних для симулятора розмірах черг
(\(n\leq16\)) асимптотика \texttt{pop}~--- \(O(n)\) через зсув масиву~---
не становить проблеми; для масштабування на більші \(N\) заплановано
перехід на циклічний буфер з \(O(1)\)-операціями. Логічний час
\texttt{simTime} збільшується методом \texttt{tickSimTime()} пропорційно
до астрономічного часу та параметра \texttt{speed}:
\(\text{simTime}_{i+1}=\text{simTime}_{i}+\Delta t_{wall}\cdot speed\),
де \(\Delta t_{wall}\)~--- реальний проміжок між кадрами. Відокремлення
логічного часу дозволяє досліджувати поведінку алгоритму у довільному
масштабі без зміни логіки активації.

Метод \texttt{execute(node, inputs)} реалізує шаблон Strategy через
перемикач за \texttt{node.kind}. \textbf{SOURCE:} вузол-джерело не має
вхідних портів; його токени вводяться ззовні через виклик
\texttt{runtime.emitFrom(nodeId, 'out', token)}, де
\(t = originT = simTime\). \textbf{DFT4:} вузол приймає 4~токени через
порти \texttt{in0}\ldots\texttt{in3} і застосовує алгоритм метелика
radix-4 з підрозд.~2.2. Реалізацію функції \texttt{computeDFT4} наведено
в лістингу~3.1.

\begin{longtable}[]{@{}
  >{\raggedright\arraybackslash}p{(\linewidth - 0\tabcolsep) * \real{0.5884}}@{}}
\caption{Реалізація 4-точкового ДПФ-метелика без нетривіальних множень.
Функція \protect\texttt{cplxMulJ(z) = \{-z.im, z.re\}} реалізує
множення на \(j\) через перестановку. Джерело: розроблено
автором.}\tabularnewline
\toprule\noalign{}
\endfirsthead
\endhead
\bottomrule\noalign{}
\endlastfoot
\begin{minipage}[t]{\linewidth}\raggedright
\begin{verbatim}
function computeDFT4(x: Complex[4]): Complex[4] {
  const E0 = cplxAdd(x[0], x[2]);   // x0 + x2
  const E1 = cplxSub(x[0], x[2]);   // x0 - x2
  const O0 = cplxAdd(x[1], x[3]);   // x1 + x3
  const O1 = cplxSub(x[1], x[3]);   // x1 - x3
  return [
    cplxAdd(E0, O0),                // X0 = E0 + O0
    cplxSub(E1, cplxMulJ(O1)),      // X1 = E1 - j*O1
    cplxSub(E0, O0),                // X2 = E0 - O0
    cplxAdd(E1, cplxMulJ(O1))       // X3 = E1 + j*O1
  ];
}
\end{verbatim}
\end{minipage} \\
\end{longtable}

Лістинг 3.1. Реалізація 4-точкового ДПФ-метелика. Джерело: розроблено автором.

\textbf{TWIDDLE:} вузол приймає 1~токен і виконує комплексне множення на
поворотний множник \(W_{16}^{k}=e^{-j2\pi k/16}\), де
\(k\in\texttt{NodeSpec.params.twiddle.k}\): \({re}'=z.re\cos\theta-z.im\sin\theta\),
\({im}'=z.re\sin\theta+z.im\cos\theta\), \(\theta=-2\pi k/16\). При
\(k=0\) вхідний токен передається без змін. \textbf{SINK:} вузол приймає
1~токен, передає його значення у \texttt{store.sinks[nodeId]} через
\texttt{JSON.stringify}, емітує \texttt{onOutput} і оновлює EMA-метрику
наскрізної затримки: \(latency=simTime-token.originT\).

Функція \texttt{generateDFT4x4(latency): Graph} (модуль
\texttt{src/graph/generateDFT4x4.ts}) реалізує шаблон Factory і будує
DFG за дев'ятьма кроками: (1)~16~вузлів \texttt{source} (\texttt{src0}
\ldots\texttt{src15}) з одним вихідним портом; (2)~4~вузли \texttt{dft4}
першого ярусу (\texttt{row-0}\ldots\texttt{row-3}) з \texttt{latency=400};
(3)~16~дуг \texttt{src}$_i\rightarrow$\texttt{row-}$r$\texttt{.in}$s$
(\(i=r+4s\)), \texttt{delay=600}; (4)~16~вузлів \texttt{twiddle}
(\texttt{tw-}$r$\texttt{-}$k_1$) з параметром \(\{N:16,\,k:r\cdot k_1\}\)
та \texttt{latency=200}; (5)~16~дуг
\texttt{row-}$r$\texttt{.out}$k_1\rightarrow$\texttt{tw-}$r$\texttt{-}$k_1$\texttt{.in};
(6)~4~вузли \texttt{dft4} другого ярусу (\texttt{col-0}\ldots\texttt{col-3});
(7)~16~дуг
\texttt{tw-}$r$\texttt{-}$k_1$\texttt{.out}$\rightarrow$\texttt{col-}$k_1$\texttt{.in}$r$;
(8)~16~вузлів \texttt{sink} (\texttt{snk0}\ldots\texttt{snk15});
(9)~16~дуг \texttt{col-}$k_1$\texttt{.out}$k_2\rightarrow$
\texttt{snk}$_{k_2+4k_1}$\texttt{.in}. Функція повертає об'єкт
\texttt{Graph} із 56~вузлами і 64~дугами. Функція
\texttt{layoutDFT4x4()} обчислює координати вузлів для React Flow:
граф розподіляється на п'ять вертикальних ярусів з кроком
\texttt{stageGap=280} пікселів; вузли в межах ярусу рівномірно
розподіляються по осі~\(y\) з кроком \texttt{rowGap=140} пікселів.

Сховище \texttt{simStore} реалізовано функцією \texttt{create()} з
Zustand і містить поля стану та методи-команди (табл.~3.3). Компоненти UI
підписуються на окремі зрізи стану через \texttt{useSimStore(selector)}
з мемоїзацією: наприклад, \texttt{Metrics} підписується лише на поле
\texttt{metrics} і перемальовується тільки при його зміні (NF3).

Таблиця 3.3.

Ключові поля та методи Zustand-сховища

\begin{longtable}[]{@{}
  >{\raggedright\arraybackslash}p{(\linewidth - 4\tabcolsep) * \real{0.2203}}
  >{\raggedright\arraybackslash}p{(\linewidth - 4\tabcolsep) * \real{0.2692}}
  >{\raggedright\arraybackslash}p{(\linewidth - 4\tabcolsep) * \real{0.4866}}@{}}
\caption[Ключові поля та методи Zustand-сховища]{Ключові поля та методи
Zustand-сховища}\tabularnewline
\toprule\noalign{}
\begin{minipage}[b]{\linewidth}\raggedright
Поле / метод
\end{minipage} & \begin{minipage}[b]{\linewidth}\raggedright
Тип
\end{minipage} & \begin{minipage}[b]{\linewidth}\raggedright
Призначення
\end{minipage} \\
\midrule\noalign{}
\endfirsthead
\toprule\noalign{}
\begin{minipage}[b]{\linewidth}\raggedright
Поле / метод
\end{minipage} & \begin{minipage}[b]{\linewidth}\raggedright
Тип
\end{minipage} & \begin{minipage}[b]{\linewidth}\raggedright
Призначення
\end{minipage} \\
\midrule\noalign{}
\endhead
\bottomrule\noalign{}
\endlastfoot
\texttt{graph} & \texttt{Graph} & Поточна топологія DFG \\
\texttt{runtime} & \texttt{DataflowRuntime|null} & Ядро симуляції \\
\texttt{running} & \texttt{boolean} & Стан rAF-циклу \\
\texttt{speed} & \texttt{number} & Множник швидкості \([0{,}1;\,10]\) \\
\texttt{simTime} & \texttt{number} & Поточний логічний час \\
\texttt{tokensByEdge} & \texttt{Record<id, Token[]>} & Токени для
анімації \\
\texttt{sinks} & \texttt{Record<id, Complex>} & Спектральні відліки
\(X(k)\) \\
\texttt{nodeActivity} & \texttt{Record<id, number>} & Теплова карта
\(a_i\) \\
\texttt{metrics} & \texttt{MetricsState} & EMA-метрики \\
\texttt{mismatches} & \texttt{Record<id, boolean>} & Результат
верифікації \\
\texttt{setGraph(g)} & \texttt{(Graph) => void} & Ініціалізація нового
графа і \texttt{runtime} \\
\texttt{injectData(v)} & \texttt{(Complex[16]) => void} & Введення
вхідного вектора \\
\texttt{start()/stop()/step()} & \texttt{() => void} & Керування
rAF-циклом \\
\texttt{applyPreset(k)} & \texttt{(PresetKind) => void} & Вбудовані
пресети вхідних сигналів \\
\end{longtable}

Компонент \texttt{TokenEdge} (модуль \texttt{src/graph/TokenEdge.tsx})
відображає кожну дугу як анімовану стрілку з кружечками-токенами. Для
кожного токена у \texttt{store.tokensByEdge[edgeId]} обчислюється
відносний прогрес руху
\(progress=(simTime-token.t_0)/token.delay\in[0,1]\) і поточна
координата лінійною інтерполяцією:
\(\mathbf{r}(p)=(1-p)\cdot\mathbf{r}_{src}+p\cdot\mathbf{r}_{tgt}\).
Колір кружечка кодує тип вузла-джерела; значення \(\{re,im\}\)
виводиться всередині кружечка. Теплова карта активності реалізується у
\texttt{simStore} через механізм експоненціального затухання з
підрозд.~2.6: функція \texttt{bumpActivity(nodeId)} викликається з
обробника \texttt{onFire} та інкрементує відповідний лічильник; у кожному
кадрі \texttt{decayActivity()} множить усі лічильники на
\(\alpha=0{,}92\).

Головний анімаційний цикл запускається у хуку \texttt{useEffect}
компонента \texttt{GraphView}. На кожному кадрі: якщо
\texttt{store.running=true}, викликається \texttt{tickSimTime()}, потім
\texttt{runtime.tick(dt, maxFires)} (у режимі \texttt{'single-fire'}
\texttt{maxFires=1}, у режимі \texttt{'run'}~--- \texttt{Infinity}); потім
\texttt{sampleQueues()} і \texttt{decayActivity()} для оновлення метрик і
теплової карти; щосекунди агрегуються лічильники \texttt{firesInWindow}
для обчислення пропускної здатності. Для уникнення блокування головного
потоку при масштабуванні на більші~\(N\) реалізовано Web Worker у файлі
\texttt{src/workers/fft.worker.ts}: через бібліотеку Comlink функція
\texttt{applyTwiddle(tokens, factors)} виконує пакетне комплексне
множення у фоновому потоці. У базовій конфігурації \(N=16\) Web Worker
не задіяний; він зарезервований для подальшого розширення.

\protect\phantomsection\label{_Toc229681146}{}3.4. Обробка виняткових
ситуацій та забезпечення надійності

Оскільки симулятор не взаємодіє з мережею, серверами чи файловою
системою, класичні джерела збоїв (мережеві тайм-аути, I/O-помилки) для
нього не характерні. Натомість існують специфічні категорії помилок,
пов'язані з валідацією вхідних даних, цілісністю графа та обмеженнями
браузерного середовища.

Під час реалізації симулятора виникли три практичні труднощі. Висока
частота оновлення інтерфейсу при великій кількості подій вимагала
чіткого розмежування обчислювального ядра і UI: стан було централізовано
у \texttt{simStore}, а компоненти переведено на селективні підписки. Для
коректної демонстрації алгоритму необхідно було відокремити логічний час
симуляції від астрономічного часу браузера; це вирішено введенням змінної
\texttt{simTime} і множника \texttt{speed}. Нарешті, для збереження
відтворюваності результатів у різних браузерах усі критичні обчислення
виконуються у фіксованому порядку обходу вузлів без мережевих викликів
у циклі симуляції.

Помилки, що можуть виникнути у симуляторі, зведено у табл.~3.4.

Таблиця 3.4.

Класифікація помилок симулятора та підходи до їх обробки

\begin{longtable}[]{@{}
  >{\raggedright\arraybackslash}p{(\linewidth - 4\tabcolsep) * \real{0.1729}}
  >{\raggedright\arraybackslash}p{(\linewidth - 4\tabcolsep) * \real{0.4115}}
  >{\raggedright\arraybackslash}p{(\linewidth - 4\tabcolsep) * \real{0.4157}}@{}}
\caption[Класифікація помилок симулятора]{Класифікація помилок симулятора
та підходи до їх обробки}\tabularnewline
\toprule\noalign{}
\begin{minipage}[b]{\linewidth}\raggedright
Категорія
\end{minipage} & \begin{minipage}[b]{\linewidth}\raggedright
Приклади
\end{minipage} & \begin{minipage}[b]{\linewidth}\raggedright
Обробка
\end{minipage} \\
\midrule\noalign{}
\endfirsthead
\toprule\noalign{}
\begin{minipage}[b]{\linewidth}\raggedright
Категорія
\end{minipage} & \begin{minipage}[b]{\linewidth}\raggedright
Приклади
\end{minipage} & \begin{minipage}[b]{\linewidth}\raggedright
Обробка
\end{minipage} \\
\midrule\noalign{}
\endhead
\bottomrule\noalign{}
\endlastfoot
Валідаційні & Введення не-числового значення у \texttt{SourceNode},
множник швидкості поза \([0{,}1;\,10]\) & Clamp до допустимого
діапазону, підсвічування червоним \\
Транзієнтні & Пропуск кадру при перевантаженні CPU & \texttt{maxFires}
обмежує кількість активацій за один \texttt{tick()} \\
Структурні & Відсутність вхідного токена на порті вузла & Вузол не
активується (правило готовності); статичний аналіз графа виключає
deadlock \\
Чисельні & Втрата точності при \(|X(k)|>10^{12}\) & Використання IEEE
754 double; попередження при виході за рамки \\
Браузерні & Відсутність Web Worker API & Graceful degradation:
обчислення у головному потоці \\
\end{longtable}

Компонент \texttt{SourceNode} використовує повзунки з межами
\([-10;\,10]\) для \texttt{re} та \texttt{im}, що виключає введення
нечислових значень. Компонент \texttt{Controls} обмежує \texttt{speed}
діапазоном \([0{,}1;\,10]\) через \texttt{Math.max(0.1, Math.min(10,
value))}. Параметри \texttt{latency} і \texttt{delay} валідуються на
додатність і цілочисельність; у разі порушення викидається виняток
\texttt{InvalidGraphError} з описовою причиною. Оскільки граф симулятора
є статичним DAG з детермінованою SDF-семантикою (підрозд.~1.2), ризик
deadlock принципово відсутній: відсутність циклів гарантує, що готовність
вузлів поступово поширюється від \texttt{source} до \texttt{sink}. Ця
властивість перевіряється на етапі побудови функцією
\texttt{validateDAG(graph)} через топологічне сортування методом Кана.

Усі значущі події ядра акумулюються у \texttt{store.eventLog} із
часовими мітками \texttt{simTime} і \texttt{performance.now()}.
Компонент \texttt{TerminalOutput} відображає журнал у реальному часі з
фільтрацією за типом події; експорт у CSV/JSON реалізований через
\texttt{Blob + URL.createObjectURL} без передачі на сервер. Рівні
логування (\texttt{debug}, \texttt{info}, \texttt{warn}, \texttt{error})
налаштовуються через \texttt{store.logLevel}; у production-збірці рівень
\texttt{debug} вимкнений за замовчуванням. При відсутності Web Workers
функція \texttt{applyTwiddle()} виконується синхронно у головному потоці
без впливу на частоту кадрів при \(N=16\). Селектори Zustand мемоїзовані
функцією \texttt{shallow}, а важкі розрахунки координат токенів на
ребрах виконуються лише для видимої області через
\texttt{ReactFlow.viewportBounds}. Середня тривалість одного кроку
\texttt{tick()} для \(N=16\) становить \(<1\)~мс, що залишає понад
15~мс запасу у бюджеті кадру (16,67~мс при 60~fps).

\protect\phantomsection\label{_Toc229681152}{}3.5. Розгортання системи

Розгортання симулятора спрощене порівняно з класичними веб-застосунками
завдяки відсутності серверної частини та бази даних. Зібраний статичний
бандл може бути розміщений на будь-якому файловому хостингу, що віддає
статичні ресурси. Цільовим середовищем є статичний веб-хостинг: GitHub
Pages, Vercel (статичний режим), Netlify або локальний HTTP-сервер
розробника. Серверна інфраструктура, бази даних, сесійне сховище та
автентифікація не застосовуються.

Збірка виконується стандартними командами Next.js (лістинг~3.2). У
результаті у директорії \texttt{out/} формується статичний бандл:
HTML-сторінка, JavaScript-бандл (Turbopack/webpack із tree-shaking і
minification) та CSS-файли; розмір gzip-бандла становить орієнтовно
250~КБ.

\begin{longtable}[]{@{}
  >{\raggedright\arraybackslash}p{(\linewidth - 0\tabcolsep) * \real{0.5586}}@{}}
\caption{Команди збірки та перевірки проєкту. Джерело: розроблено
автором.}\tabularnewline
\toprule\noalign{}
\endfirsthead
\endhead
\bottomrule\noalign{}
\endlastfoot
\begin{minipage}[t]{\linewidth}\raggedright
\begin{verbatim}
$ pnpm install            # встановлення залежностей
$ pnpm run build          # next build -> .next/
$ pnpm run export         # next export -> out/
$ pnpm run lint           # ESLint
$ pnpm test               # Jest-тести
\end{verbatim}
\end{minipage} \\
\end{longtable}

Лістинг 3.2. Команди збірки та перевірки проєкту. Джерело: розроблено автором.

Для автоматизації перевірок налаштовано GitHub Actions із workflow
\texttt{.github/workflows/ci.yml}: job \texttt{lint} запускає ESLint і
TypeScript-компіляцію у режимі \texttt{--noEmit}; job \texttt{test}
запускає Jest з генерацією звіту про покриття; job \texttt{build}
перевіряє, що збірка проходить без помилок. Стратегія релізу ---
\emph{rolling}: нова статична збірка замінює попередню при успішному
проходженні CI. Оскільки застосунок не має сесійного стану, жодних
процедур міграції бази даних чи інвалідації кешу не потрібно.
Контейнеризація, оркестрація та резервне копіювання бази даних для цього
проєкту не застосовуються, оскільки система розгортається як статичний
клієнтський застосунок без серверної логіки.

Оскільки вихідний код є єдиним артефактом, що підлягає збереженню,
резервне копіювання реалізоване через Git: повна історія зберігається у
віддаленому репозиторії GitHub. Детермінованість збірки гарантується
файлом \texttt{pnpm-lock.yaml}. Будь-який виконавець може відтворити
продуктову збірку командою \texttt{pnpm install \&\& pnpm run build} на
машині з Node.js~\(\geq18\).

\protect\phantomsection\label{_Toc229681157}{}3.6. Тестування програмного
забезпечення

Тестування на етапі розробки охоплює модульні тести найкритичніших
компонентів і базові інтеграційні сценарії. Комплексне тестування готової
системи і аналіз продуктивності розглядаються у розділі~4.

Для модульного тестування використовується \textbf{Jest~29.x} з
конфігурацією у файлі \texttt{jest.config.ts}: середовище \texttt{node}
для модулів \texttt{core/} і \texttt{utils/} та \texttt{jsdom} для
UI-тестів; трансформатор \texttt{ts-jest} з \texttt{tsconfig.test.json};
поріг покриття --- не менше 70~\% тверджень, віток і функцій у
\texttt{src/core/} та \texttt{src/utils/}.

Реалізовані тест-сьюти та цільові модулі наведено у табл.~3.5.
Лістинг~3.3 наводить приклад тесту коректності функції
\texttt{computeDFT4} для імпульсного входу \((1,0,0,0)\): очікуваний
результат --- \((1,1,1,1)\), оскільки всі поворотні множники у сумі
дають одиничні коефіцієнти.

Таблиця 3.5.

Модульні тести симулятора

\begin{longtable}[]{@{}
  >{\raggedright\arraybackslash}p{(\linewidth - 4\tabcolsep) * \real{0.1885}}
  >{\raggedright\arraybackslash}p{(\linewidth - 4\tabcolsep) * \real{0.2904}}
  >{\raggedright\arraybackslash}p{(\linewidth - 4\tabcolsep) * \real{0.5212}}@{}}
\caption[Модульні тести симулятора]{Модульні тести симулятора}\tabularnewline
\toprule\noalign{}
\begin{minipage}[b]{\linewidth}\raggedright
Тест-файл
\end{minipage} & \begin{minipage}[b]{\linewidth}\raggedright
Цільовий модуль
\end{minipage} & \begin{minipage}[b]{\linewidth}\raggedright
Сценарії
\end{minipage} \\
\midrule\noalign{}
\endfirsthead
\toprule\noalign{}
\begin{minipage}[b]{\linewidth}\raggedright
Тест-файл
\end{minipage} & \begin{minipage}[b]{\linewidth}\raggedright
Цільовий модуль
\end{minipage} & \begin{minipage}[b]{\linewidth}\raggedright
Сценарії
\end{minipage} \\
\midrule\noalign{}
\endhead
\bottomrule\noalign{}
\endlastfoot
\texttt{math.test.ts} & \texttt{utils/math.ts} & Коректність
\texttt{cplxAdd}, \texttt{cplxSub}, \texttt{cplxMulJ}; правильність
\texttt{computeDFT4} на референтних входах \\
\texttt{buffers.test.ts} & \texttt{core/buffers.ts} & FIFO-порядок,
поведінка \texttt{peek} на пустій черзі, розмір \\
\texttt{scheduler.test.ts} & \texttt{core/scheduler.ts} & Правило
активації, коректність \texttt{tick()}, дотримання \texttt{maxFires},
детермінованість \\
\texttt{generate.test.ts} & \texttt{graph/generateDFT4x4.ts} & Кількість
вузлів (56) і дуг (64), коректність з'єднань \\
\texttt{compare.test.ts} & \texttt{utils/compare.ts} & Збіг симулятора з
fft.js на імпульсі та синусоїді з \(\varepsilon\leq10^{-9}\) \\
\texttt{store.test.ts} & \texttt{store/simStore.ts} & Скидання стану,
\texttt{injectData}, \texttt{applyPreset} \\
\end{longtable}

\begin{longtable}[]{@{}
  >{\raggedright\arraybackslash}p{(\linewidth - 0\tabcolsep) * \real{0.5817}}@{}}
\caption{Приклад модульного тесту функції \protect\texttt{computeDFT4}.
Джерело: розроблено автором.}\tabularnewline
\toprule\noalign{}
\endfirsthead
\endhead
\bottomrule\noalign{}
\endlastfoot
\begin{minipage}[t]{\linewidth}\raggedright
\begin{verbatim}
describe('computeDFT4', () => {
  it('returns (1,1,1,1) for unit impulse (1,0,0,0)', () => {
    const input: Complex[] = [
      { re: 1, im: 0 }, { re: 0, im: 0 },
      { re: 0, im: 0 }, { re: 0, im: 0 }
    ];
    const out = computeDFT4(input);
    for (const z of out) {
      expect(z.re).toBeCloseTo(1, 12);
      expect(z.im).toBeCloseTo(0, 12);
    }
  });
});
\end{verbatim}
\end{minipage} \\
\end{longtable}

Лістинг 3.3. Приклад модульного тесту функції \texttt{computeDFT4}. Джерело: розроблено автором.

Поточний рівень покриття, отриманий командою \texttt{pnpm test
--coverage}, становить: \texttt{src/core/}~--- 85~\% statements,
78~\% branches, 92~\% functions; \texttt{src/utils/}~--- 91~\%
statements, 84~\% branches, 100~\% functions;
\texttt{src/graph/generateDFT4x4.ts}~--- 100~\% statements. Усі
значення перевищують цільовий поріг 70~\%. UI-компоненти
(\texttt{src/components/}) охоплені лише дим-тестами на монтування і
не враховуються у покритті, оскільки їх коректність підтверджується
ручним тестуванням у розділі~4. Тести автоматично виконуються у GitHub
Actions на кожному push до \texttt{main}; звіт у форматі HTML
публікується як артефакт збірки.

\protect\phantomsection\label{_Toc229681163}{}3.7. Висновки до розділу~3

У третьому розділі подано програмну реалізацію симулятора, що відповідає
архітектурі, спроєктованій у розділі~2.

Описано технологічний стек: TypeScript~5 як статично типізована мова,
Next.js~14 як фреймворк статичної SPA, React~18 і React Flow~12 для
реактивного UI та графової візуалізації, Zustand~4 для сховища стану,
Jest~29 для модульних тестів, ESLint/Prettier для забезпечення стилю.
Усі бібліотеки мають відкриті ліцензії MIT або Apache~2.0 (NF6).

Організовано вихідний код за принципом розподілу за шарами
відповідальності: ядро симуляції (\texttt{src/core/}) без залежностей від
UI, генерація та візуалізація графа (\texttt{src/graph/}),
UI-компоненти (\texttt{src/components/}), сховище стану
(\texttt{src/store/}), утиліти (\texttt{src/utils/}) та Web Workers
(\texttt{src/workers/}). Повну відповідність архітектурних компонентів
до модулів коду подано у табл.~3.2.

Реалізовано ключові алгоритми: клас \texttt{DataflowRuntime} з методом
\texttt{tick()} у чотири кроки (доставка \(\rightarrow\) перевірка
готовності \(\rightarrow\) активація \(\rightarrow\) розміщення вихідних
токенів); функція \texttt{generateDFT4x4()} як Factory для побудови DFG
із 56~вузлів і 64~дуг; сховище \texttt{simStore} на Zustand із
підпискою UI на окремі зрізи стану; компоненти \texttt{TokenEdge} та
теплової карти з інтерполяцією і експоненціальним затуханням активності.

Розглянуто обробку виняткових ситуацій: класифіковано помилки за
категоріями (валідаційні, транзієнтні, структурні, чисельні, браузерні),
доведено відсутність ризику deadlock у DAG-структурі, описано механізми
логування, graceful degradation та оптимізації перерендерингу. Визначено
процедуру розгортання як статичної SPA на довільному файловому хостингу;
налаштовано GitHub Actions для автоматичної перевірки типів, лінтера та
тестів. Проведено модульне тестування у Jest з покриттям 85~\% для
\texttt{src/core/} та 91~\% для \texttt{src/utils/}, що перевищує
цільовий поріг 70~\%; реалізовано шість тест-сьютів, що охоплюють
комплексну арифметику, FIFO-черги, планувальник, генератор графа,
верифікацію та сховище стану. Комплексна експериментальна оцінка якості
--- точність відносно еталонного fft.js, продуктивність, аналіз
критичного шляху --- виноситься до розділу~4.